\begin{table}
\centering
\caption{Effect of Pirate Encounters on Trip Frequency. \label{tab:trip-count}}
\centering
\begin{adjustbox}{width = .9\textwidth}
\begin{threeparttable}
\begin{tabular}[t]{lcccc}
\toprule
  & Global & G. of Aden & G. of Guinea & S.E. Asia\\
\midrule
Avg. Encounters (7-day lag) & 0.0054* & 0.0156* & -0.0050 & -0.0006\\
 & (0.0032) & (0.0088) & (0.0039) & (0.0041)\\
 &  &  &  & \\
Observations & 968,785 & 98,564 & 115,310 & 114,351\\
 &  &  &  & \\
Country-Pair FE & X & X & X & X\\
Hotspot FE & X & $\bullet$ & $\bullet$ & $\bullet$\\
Month-by-Year FE & X & X & X & X\\
\bottomrule
\end{tabular}
\begin{tablenotes}
\item \scriptsize * p $<$ 0.1, ** p $<$ 0.05, *** p $<$ 0.01 The unit of observation is a country-pair by 7-day window. The outcome is the count of cargo vessel trips on that route in that week. The model is Negative Binomial; coefficients are log-multiplicative and approximately equal to percentage changes in trips for small values. The sample spans from 2012 to 2023. Global uses the full sample; G. of Aden, G. of Guinea, and S.E. Asia restrict the sample to routes passing through each hotspot, respectively. Avg. Encounters (7-day lag) is the average of each voyage's 7-day pirate encounter count in the prior 7-day window on the same route, using a 5-degree spatial footprint. This mirrors the 7-day encounter window used in the main voyage-level analysis. Controls include average wind speed, the wind-resistance index, and average wave height. Fixed effects include country-to-country combination, hotspot, and month-by-year. X indicates the fixed effect is included; $\bullet$ indicates the fixed effect is absorbed by the sample restriction. Standard errors are clustered by country-to-country route by year.
\end{tablenotes}
\end{threeparttable}
\end{adjustbox}
\end{table}
